\section{Introduction}
\label{sec:intro}

Modern OFDM and MIMO wireless systems rely on accurate channel state information (CSI) for tasks such as equalization, precoding, and scheduling. CSI is typically obtained by transmitting known pilot symbols on a subset of subcarriers and estimating the channel from the resulting observations. Every pilot takes up time--frequency resources that could otherwise carry data, and in fast-changing or latency-sensitive settings, the time spent collecting pilots directly delays when a usable channel estimate becomes available.

A common approach is to fix the pilot pattern in advance (e.g., uniformly scattered placement) and then estimate the channel with a batch estimator, such as the linear minimum mean-square error (LMMSE) method, once all pilots have been collected~\cite{yu2025lmmse}. This approach is simple and well understood, but it commits to a pilot budget and to specific pilot locations ahead of time, without taking into account what has already been learned from earlier measurements.

Adaptive and compressive pilot allocation schemes have been proposed to reduce this overhead by exploiting channel sparsity or by adjusting pilot density based on channel conditions~\cite{arai2025joint}. Related ideas also appear in sequential and active sensing, where measurements are chosen one at a time to reduce uncertainty as quickly as possible~\cite{sohrabi2021active}. However, these methods are not usually framed as a simple, fully model-based rule that updates an LMMSE estimate online and decides, one subcarrier at a time, which measurement is most worth taking next.

In this work, we propose an incremental pilot-selection framework that measures a single subcarrier at a time, updates an LMMSE posterior after each measurement, and selects the next subcarrier by greedily maximizing the predicted reduction in the posterior error-covariance trace. The method is entirely model-based and requires no training. The rest of this paper is organized as follows: Section~\ref{sec:system_model} describes the system model and problem formulation, Section~\ref{sec:method} details the proposed active selection rule, Section~\ref{sec:experiments} presents the experimental study, and Section~\ref{sec:conclusion} concludes the paper.
